% ---- score2 standard preamble (from: python3 score2_tools.py preamble) ----
\documentclass[10pt]{extarticle}
\usepackage[a4paper,margin=1.5cm]{geometry}
\usepackage[T1]{fontenc}
\usepackage{amsmath,amssymb}
\usepackage{graphicx}
\usepackage[export]{adjustbox}
\usepackage{multicol}
\usepackage{enumitem}
\usepackage{booktabs,array,tabularx}
\setlength{\columnsep}{0.7cm}
\setlength{\parindent}{0pt}
\setlength{\parskip}{3pt}
\pagestyle{empty}
\makeatletter
% figure registry (figdefs.tex is generated by the extract step)
\newcommand{\qfigdef}[4]{\@namedef{qf@file@#1}{#2}\@namedef{qf@w@#1}{#3}\@namedef{qf@wide@#1}{#4}}
\newcommand{\qf@check}[1]{\@ifundefined{qf@file@#1}{\PackageError{score2}{Unknown figure id #1}{}}{}}
\newcommand{\qf@img}[4][]{\includegraphics[width=#2\dimexpr\@nameuse{qf@w@#4}pt\relax,max width=#3,max height=0.45\textheight,#1]{\@nameuse{qf@file@#4}}}
% Full-width questions: finalize wraps every question that contains a wide figure in \qfullbegin...\qfullend,
% so the WHOLE question (stem, figures, options) is set across the page in reading order, then the two columns
% resume. The question's head - its number, stem and first figure (\qfullsplit marks the end of the head) -
% is kept on one page, so a number is never left at a page foot with its figure overleaf; the rest may break
% across pages normally, so no big gaps are left. Without a \qfullsplit, the whole question is kept together
% if it is under 45% of a page. Never write these yourself.
\newif\ifqf@full
\newif\ifqf@boxed
\newsavebox{\qf@box}
\newcommand{\qf@open}{\global\qf@fulltrue\par\global\qf@boxedtrue
  \setbox\qf@box\vbox\bgroup\hsize\textwidth\linewidth\textwidth}
\newcommand{\qf@emit}[1]{\par\egroup\global\qf@boxedfalse
  \ifdim\dimexpr\ht\qf@box+\dp\qf@box\relax<#1\textheight\noindent\usebox{\qf@box}\par\else\unvbox\qf@box\fi}
\newcommand{\qfullsplit}{\ifqf@boxed\qf@emit{0.85}\fi}
\newcommand{\qf@close}{\ifqf@boxed\qf@emit{0.45}\fi\global\qf@fullfalse}
\newcommand{\qfullbegin}{\end{multicols}\qf@open}
\newcommand{\qfullend}{\qf@close\begin{multicols}{2}}
\newcommand{\qfullstart}{\qf@open}
\newcommand{\qfulldone}{\qf@close}
% \qfig[scale]{id}: block figure, centred, directly where it is written (size set by the tool).
% Fallback if a wide figure is not inside a full-width question: break out of the columns in place - no floats.
\newcommand{\qfig}[2][1]{\qf@check{#2}%
  \ifqf@full\qf@blk{#1}{#2}\else
  \ifnum\@nameuse{qf@wide@#2}=1 \expandafter\@firstoftwo\else\expandafter\@secondoftwo\fi
  {\end{multicols}\begin{center}\qf@img{#1}{\textwidth}{#2}\end{center}\begin{multicols}{2}}%
  {\qf@blk{#1}{#2}}\fi}
\newcommand{\qf@blk}[2]{\par\smallskip{\centering\qf@img{#1}{\linewidth}{#2}\par}\smallskip}
% \qfigi[scale]{id}: inline figure (inside option lists / tables / minipages); never breaks columns.
\newcommand{\qfigi}[2][1]{\qf@check{#2}\qf@img[valign=c]{#1}{\linewidth}{#2}}
\makeatother
% structure macros - lint reads these to check completeness and figure placement
\newcommand{\papersec}[1]{\par\bigskip{\centering\Large\bfseries PAPER-#1\par}\medskip}
\newcommand{\partsec}[2]{\par\medskip\noindent{\large\bfseries PART-#1: #2}\par\smallskip}
\newcommand{\sect}[1]{\par\smallskip\noindent\textbf{#1}\par}
\newcommand{\qn}[1]{\par\smallskip\noindent\textbf{#1.}~}
\newcommand{\paraq}[2]{\par\smallskip\noindent\textbf{Paragraph for Questions #1 and #2}\par}
\newenvironment{qoptions}{\begin{itemize}[leftmargin=1.8em,itemsep=1pt,topsep=2pt,label={}]}{\end{itemize}}
\newcommand{\opt}[2]{\item[(#1)] #2}
% figure registry (generated by score2_tools.py finalize)
\qfigdef{p1_phy_s1_q01_1}{TOAS - Test 05 (03 Apr 2022)-img/p1_phy_s1_q01_1.png}{97.9}{0}
\qfigdef{p1_phy_s1_q02_1}{TOAS - Test 05 (03 Apr 2022)-img/p1_phy_s1_q02_1.png}{109.9}{0}
\qfigdef{p1_phy_s1_q03_1}{TOAS - Test 05 (03 Apr 2022)-img/p1_phy_s1_q03_1.png}{143.7}{0}
\qfigdef{p1_phy_s1_q04_1}{TOAS - Test 05 (03 Apr 2022)-img/p1_phy_s1_q04_1.png}{152.0}{0}
\qfigdef{p1_phy_s1_q05_1}{TOAS - Test 05 (03 Apr 2022)-img/p1_phy_s1_q05_1.png}{130.1}{0}
\qfigdef{p1_phy_s1_q06_1}{TOAS - Test 05 (03 Apr 2022)-img/p1_phy_s1_q06_1.png}{95.0}{0}
\qfigdef{p1_phy_s1_q07_1}{TOAS - Test 05 (03 Apr 2022)-img/p1_phy_s1_q07_1.png}{137.3}{0}
\qfigdef{p1_phy_s1_q08_1}{TOAS - Test 05 (03 Apr 2022)-img/p1_phy_s1_q08_1.png}{128.8}{0}
\qfigdef{p1_phy_s2_q11_1}{TOAS - Test 05 (03 Apr 2022)-img/p1_phy_s2_q11_1.png}{80.8}{0}
\qfigdef{p1_phy_s2_q12_1}{TOAS - Test 05 (03 Apr 2022)-img/p1_phy_s2_q12_1.png}{90.8}{0}
\qfigdef{p1_phy_s3_q01_1}{TOAS - Test 05 (03 Apr 2022)-img/p1_phy_s3_q01_1.png}{78.1}{0}
\qfigdef{p1_phy_s3_q02_1}{TOAS - Test 05 (03 Apr 2022)-img/p1_phy_s3_q02_1.png}{81.6}{0}
\qfigdef{p1_phy_s3_q03_1}{TOAS - Test 05 (03 Apr 2022)-img/p1_phy_s3_q03_1.png}{102.7}{0}
\qfigdef{p1_phy_s3_q04_1}{TOAS - Test 05 (03 Apr 2022)-img/p1_phy_s3_q04_1.png}{143.2}{0}
\qfigdef{p1_phy_s3_q05_1}{TOAS - Test 05 (03 Apr 2022)-img/p1_phy_s3_q05_1.png}{140.6}{0}
\qfigdef{p1_chem_s1_q01_A}{TOAS - Test 05 (03 Apr 2022)-img/p1_chem_s1_q01_A.png}{131.8}{0}
\qfigdef{p1_chem_s1_q01_B}{TOAS - Test 05 (03 Apr 2022)-img/p1_chem_s1_q01_B.png}{139.6}{0}
\qfigdef{p1_chem_s1_q01_C}{TOAS - Test 05 (03 Apr 2022)-img/p1_chem_s1_q01_C.png}{143.0}{0}
\qfigdef{p1_chem_s1_q01_D}{TOAS - Test 05 (03 Apr 2022)-img/p1_chem_s1_q01_D.png}{146.2}{0}
\qfigdef{p1_chem_s1_q02_A}{TOAS - Test 05 (03 Apr 2022)-img/p1_chem_s1_q02_A.png}{90.6}{0}
\qfigdef{p1_chem_s1_q02_B}{TOAS - Test 05 (03 Apr 2022)-img/p1_chem_s1_q02_B.png}{78.5}{0}
\qfigdef{p1_chem_s1_q02_C}{TOAS - Test 05 (03 Apr 2022)-img/p1_chem_s1_q02_C.png}{90.3}{0}
\qfigdef{p1_chem_s1_q02_D}{TOAS - Test 05 (03 Apr 2022)-img/p1_chem_s1_q02_D.png}{29.7}{0}
\qfigdef{p1_chem_s1_q03_1}{TOAS - Test 05 (03 Apr 2022)-img/p1_chem_s1_q03_1.png}{58.0}{0}
\qfigdef{p1_chem_s3_q01_1}{TOAS - Test 05 (03 Apr 2022)-img/p1_chem_s3_q01_1.png}{293.1}{1}
\qfigdef{p1_chem_s3_q02_2}{TOAS - Test 05 (03 Apr 2022)-img/p1_chem_s3_q02_2.png}{178.8}{0}
\qfigdef{p1_chem_s3_q03_1}{TOAS - Test 05 (03 Apr 2022)-img/p1_chem_s3_q03_1.png}{243.5}{0}
\qfigdef{p1_math_s1_q03_1}{TOAS - Test 05 (03 Apr 2022)-img/p1_math_s1_q03_1.png}{84.0}{0}
\qfigdef{p1_math_s1_q07_1}{TOAS - Test 05 (03 Apr 2022)-img/p1_math_s1_q07_1.png}{105.3}{0}
\qfigdef{p2_phy_s1_q03_1}{TOAS - Test 05 (03 Apr 2022)-img/p2_phy_s1_q03_1.png}{128.0}{0}
\qfigdef{p2_phy_s1_q05_1}{TOAS - Test 05 (03 Apr 2022)-img/p2_phy_s1_q05_1.png}{112.7}{0}
\qfigdef{p2_phy_s1_q06_1}{TOAS - Test 05 (03 Apr 2022)-img/p2_phy_s1_q06_1.png}{104.3}{0}
\qfigdef{p2_phy_s1_q08_1}{TOAS - Test 05 (03 Apr 2022)-img/p2_phy_s1_q08_1.png}{33.7}{0}
\qfigdef{p2_phy_s1_q09_1}{TOAS - Test 05 (03 Apr 2022)-img/p2_phy_s1_q09_1.png}{137.9}{0}
\qfigdef{p2_phy_s1_q10_1}{TOAS - Test 05 (03 Apr 2022)-img/p2_phy_s1_q10_1.png}{84.2}{0}
\qfigdef{p2_phy_s1_q10_A}{TOAS - Test 05 (03 Apr 2022)-img/p2_phy_s1_q10_A.png}{69.3}{0}
\qfigdef{p2_phy_s1_q10_B}{TOAS - Test 05 (03 Apr 2022)-img/p2_phy_s1_q10_B.png}{74.0}{0}
\qfigdef{p2_phy_s1_q10_C}{TOAS - Test 05 (03 Apr 2022)-img/p2_phy_s1_q10_C.png}{79.9}{0}
\qfigdef{p2_phy_s1_q10_D}{TOAS - Test 05 (03 Apr 2022)-img/p2_phy_s1_q10_D.png}{75.2}{0}
\qfigdef{p2_phy_s1_q11_1}{TOAS - Test 05 (03 Apr 2022)-img/p2_phy_s1_q11_1.png}{70.5}{0}
\qfigdef{p2_phy_s1_q16_A}{TOAS - Test 05 (03 Apr 2022)-img/p2_phy_s1_q16_A.png}{104.1}{0}
\qfigdef{p2_phy_s1_q16_B}{TOAS - Test 05 (03 Apr 2022)-img/p2_phy_s1_q16_B.png}{106.0}{0}
\qfigdef{p2_phy_s1_q16_C}{TOAS - Test 05 (03 Apr 2022)-img/p2_phy_s1_q16_C.png}{102.6}{0}
\qfigdef{p2_phy_s1_q16_D}{TOAS - Test 05 (03 Apr 2022)-img/p2_phy_s1_q16_D.png}{106.0}{0}
\qfigdef{p2_phy_s1_q17_1}{TOAS - Test 05 (03 Apr 2022)-img/p2_phy_s1_q17_1.png}{102.4}{0}
\qfigdef{p2_phy_s1_q18_1}{TOAS - Test 05 (03 Apr 2022)-img/p2_phy_s1_q18_1.png}{93.0}{0}
\qfigdef{p2_phy_s1_q19_1}{TOAS - Test 05 (03 Apr 2022)-img/p2_phy_s1_q19_1.png}{103.1}{0}
\qfigdef{p2_phy_s2_q02_1}{TOAS - Test 05 (03 Apr 2022)-img/p2_phy_s2_q02_1.png}{147.0}{0}
\qfigdef{p2_phy_s2_q03_1}{TOAS - Test 05 (03 Apr 2022)-img/p2_phy_s2_q03_1.png}{131.1}{0}
\qfigdef{p2_phy_s2_q04_1}{TOAS - Test 05 (03 Apr 2022)-img/p2_phy_s2_q04_1.png}{63.0}{0}
\qfigdef{p2_phy_s2_q05_1}{TOAS - Test 05 (03 Apr 2022)-img/p2_phy_s2_q05_1.png}{101.1}{0}
\qfigdef{p2_phy_s2_q06_1}{TOAS - Test 05 (03 Apr 2022)-img/p2_phy_s2_q06_1.png}{72.8}{0}
\qfigdef{p2_phy_s2_q09_1}{TOAS - Test 05 (03 Apr 2022)-img/p2_phy_s2_q09_1.png}{120.8}{0}
\qfigdef{p2_phy_s2_q10_1}{TOAS - Test 05 (03 Apr 2022)-img/p2_phy_s2_q10_1.png}{85.8}{0}
\qfigdef{p2_chem_s1_q01_1}{TOAS - Test 05 (03 Apr 2022)-img/p2_chem_s1_q01_1.png}{127.5}{0}
\qfigdef{p2_chem_s1_q02_1}{TOAS - Test 05 (03 Apr 2022)-img/p2_chem_s1_q02_1.png}{81.1}{0}
\qfigdef{p2_chem_s1_q02_A}{TOAS - Test 05 (03 Apr 2022)-img/p2_chem_s1_q02_A.png}{45.5}{0}
\qfigdef{p2_chem_s1_q02_B}{TOAS - Test 05 (03 Apr 2022)-img/p2_chem_s1_q02_B.png}{58.3}{0}
\qfigdef{p2_chem_s1_q02_C}{TOAS - Test 05 (03 Apr 2022)-img/p2_chem_s1_q02_C.png}{46.3}{0}
\qfigdef{p2_chem_s1_q02_D}{TOAS - Test 05 (03 Apr 2022)-img/p2_chem_s1_q02_D.png}{53.4}{0}
\qfigdef{p2_chem_s1_q03_1}{TOAS - Test 05 (03 Apr 2022)-img/p2_chem_s1_q03_1.png}{120.5}{0}
\qfigdef{p2_chem_s1_q03_A}{TOAS - Test 05 (03 Apr 2022)-img/p2_chem_s1_q03_A.png}{25.0}{0}
\qfigdef{p2_chem_s1_q03_B}{TOAS - Test 05 (03 Apr 2022)-img/p2_chem_s1_q03_B.png}{25.2}{0}
\qfigdef{p2_chem_s1_q03_C}{TOAS - Test 05 (03 Apr 2022)-img/p2_chem_s1_q03_C.png}{25.5}{0}
\qfigdef{p2_chem_s1_q04_1}{TOAS - Test 05 (03 Apr 2022)-img/p2_chem_s1_q04_1.png}{141.1}{0}
\qfigdef{p2_chem_s1_q04_A}{TOAS - Test 05 (03 Apr 2022)-img/p2_chem_s1_q04_A.png}{55.1}{0}
\qfigdef{p2_chem_s1_q04_B}{TOAS - Test 05 (03 Apr 2022)-img/p2_chem_s1_q04_B.png}{55.1}{0}
\qfigdef{p2_chem_s1_q04_C}{TOAS - Test 05 (03 Apr 2022)-img/p2_chem_s1_q04_C.png}{52.9}{0}
\qfigdef{p2_chem_s1_q04_D}{TOAS - Test 05 (03 Apr 2022)-img/p2_chem_s1_q04_D.png}{52.9}{0}
\qfigdef{p2_chem_s1_q05_A}{TOAS - Test 05 (03 Apr 2022)-img/p2_chem_s1_q05_A.png}{40.7}{0}
\qfigdef{p2_chem_s1_q05_B}{TOAS - Test 05 (03 Apr 2022)-img/p2_chem_s1_q05_B.png}{30.4}{0}
\qfigdef{p2_chem_s1_q05_C}{TOAS - Test 05 (03 Apr 2022)-img/p2_chem_s1_q05_C.png}{43.0}{0}
\qfigdef{p2_chem_s1_q05_D}{TOAS - Test 05 (03 Apr 2022)-img/p2_chem_s1_q05_D.png}{41.9}{0}
\qfigdef{p2_chem_s1_q06_1}{TOAS - Test 05 (03 Apr 2022)-img/p2_chem_s1_q06_1.png}{75.6}{0}
\qfigdef{p2_chem_s1_q06_2}{TOAS - Test 05 (03 Apr 2022)-img/p2_chem_s1_q06_2.png}{129.3}{0}
\qfigdef{p2_chem_s1_q20_1}{TOAS - Test 05 (03 Apr 2022)-img/p2_chem_s1_q20_1.png}{82.0}{0}
\qfigdef{p2_chem_s2_q01_1}{TOAS - Test 05 (03 Apr 2022)-img/p2_chem_s2_q01_1.png}{123.4}{0}
\qfigdef{p2_chem_s2_q02_1}{TOAS - Test 05 (03 Apr 2022)-img/p2_chem_s2_q02_1.png}{37.5}{0}
\qfigdef{p2_chem_s2_q02_2}{TOAS - Test 05 (03 Apr 2022)-img/p2_chem_s2_q02_2.png}{113.1}{0}
\qfigdef{p2_chem_s2_q03_1}{TOAS - Test 05 (03 Apr 2022)-img/p2_chem_s2_q03_1.png}{59.0}{0}
\qfigdef{p2_math_s1_q20_1}{TOAS - Test 05 (03 Apr 2022)-img/p2_math_s1_q20_1.png}{74.9}{0}
\raggedcolumns
\makeatletter\@beginparpenalty=10000\makeatother
\newcolumntype{L}[1]{>{\raggedright\arraybackslash}p{#1}}
\begin{document}
\begin{center}{\Large\bfseries ALLEN TOAS --- Test 05 (03 Apr 2022)}\end{center}
\begin{multicols}{2}
\raggedcolumns
\papersec{1}
\partsec{1}{PHYSICS}
\sect{SECTION-I(i)}
\qn{1} Table shown in figure is rotating about its vertical axis with an angular speed $\omega=2$ rad/sec. Blocks $A$ and $B$ have masses 45 kg and 22.5 kg respectively. String connecting them is light and passes through a small hole in the table. Pulley shown in massless. Coefficient of friction at the contacts of $A$ and $B$ with the table is 0.4. dimensions of two blocks are small and distance of $B$ from rotation axis is 0.6 m. If block $A$ remains stationary relative to the table at a distance `$x$' from the rotation axis then (Take : $g=10$ m/s$^2$)
\qfig{p1_phy_s1_q01_1}
\begin{qoptions}
\opt{A}{Minimum possible value of $x$ is 0.13 m}
\opt{B}{Minimum possible value of $x$ is 0.26 m}
\opt{C}{Maximum possible value of $x$ is 2.37 m}
\opt{D}{Maximum possible value of $x$ is 1.37 m}
\end{qoptions}
\qn{2} A ball of mass $M$ is swing around in a circle in a horizontal plane. On end of spring (spring constant $k$) is fixed and other end is connected to the ball as shown in figure. The ball describes a horizontal circle a distance $h$ above the horizontal floor with constant speed. The stretched spring has a length $\ell$ and makes an angle $60^\circ$ with the vertical. Neglect air resistance.
\qfig{p1_phy_s1_q02_1}
\begin{qoptions}
\opt{A}{Magnitude of force exerted by spring on block is $2\,Mg$}
\opt{B}{Natural length of spring is $\left(\ell+\dfrac{2Mg}{k}\right)$.}
\opt{C}{Speed of the ball is $\sqrt{\dfrac{3g\ell}{2}}$}
\opt{D}{At some instant of time, if spring breaks, the ball moves a horizontal distance $\sqrt{3h\ell}$, before it hits the floor.}
\end{qoptions}
\qn{3} A small toy car of mass $m$ slides with negligible friction on a loop track as shown in figure. The toy car starts from rest at a point $H=4R$ above the lowest point of the track. (Acceleration due to gravity is $g$ and radius of circular part is $R$).
\qfig{p1_phy_s1_q03_1}
\begin{qoptions}
\opt{A}{Speed of toy car at point $B$ is $\sqrt{4gR}$}
\opt{B}{Normal force exerted by track on the toy car at point $B$ is 9 $mg$.}
\opt{C}{Normal force exerted by track on the toy car at position $C$ is 6 $mg$.}
\opt{D}{Toy car looses contact with track at a point between $C$ and $D$.}
\end{qoptions}
\qn{4} A small block slides down a hill (starting from rest) of height $H$. The part of the hill near its end $B$ is horizontal. The block flies off the hill at $B$ to hit the ground at $C$. Hill is smooth. If $y$ is height of end $B$ of hill from ground and $x$ is horizontal displacement of block before hitting the ground after loosing contact with track, then :
\qfig{p1_phy_s1_q04_1}
\begin{qoptions}
\opt{A}{For $x$ to be maximum, $y$ is $H$}
\opt{B}{For $x$ to be maximum $y$ is $\dfrac H2$.}
\opt{C}{Maximum possible value of $x$ is $\dfrac H2$.}
\opt{D}{Maximum possible value of $x$ is $H$.}
\end{qoptions}
\qn{5} A thin non conducting ring of radius $R$ has linear charge density $\lambda$ varying as $\lambda=\lambda_0\cos\theta$, where $\lambda_0$ is a positive constant and $\theta$ is angle measured from positive $x$-direction as shown. Consider centre of ring is at origin.
\qfig{p1_phy_s1_q05_1}
\begin{qoptions}
\opt{A}{Resultant electric field at origin is towards positive $y$-axis.}
\opt{B}{Resultant electric field at origin is towards negative $x$-axis.}
\opt{C}{Magnitude of resultant electric field at origin is $\dfrac{\lambda_0}{8\varepsilon_0R}$.}
\opt{D}{Magnitude of resultant electric field at origin is $\dfrac{\lambda_0}{4\varepsilon_0R}$.}
\end{qoptions}
\qn{6} A fixed smooth inclined plane making an angle $\theta$ with the horizontal is placed in a uniform horizontal electric field $E_0$ as shown in figure. A small block of mass $m$ having charge $+q$ is released on inclined plane. (Acceleration due to gravity is `$g$')
\qfig{p1_phy_s1_q06_1}
\begin{qoptions}
\opt{A}{Normal contact force applied by inclined plane on block is $(mg\cos\theta+qE_0\sin\theta)$}
\opt{B}{Normal contact force force applied by inclined plane on block is $(mg\cos\theta-qE_0\sin\theta)$}
\opt{C}{If acceleration of block is zero then $\theta=\tan^{-1}\left(\dfrac{qE_0}{mg}\right)$.}
\opt{D}{If acceleration of block is zero then $\theta=\tan^{-1}\left(\dfrac{mg}{qE_0}\right)$.}
\end{qoptions}
\qn{7} A long cylindrical tube containing water is closed by an equiconvex lens of focal length 10 cm in air. A point object `$O$' is placed along the axis of the tube outside it at a distance of 21 cm from the lens. Refractive index of material of lens is 1.5 and that of water is $\dfrac43$.
\qfig{p1_phy_s1_q07_1}
\begin{qoptions}
\opt{A}{Final image is virtual in nature.}
\opt{B}{Final image is real in nature.}
\opt{C}{Final image is located at a distance 70 cm from lens inside the tube.}
\opt{D}{Final image is located at a distance 30 cm from lens inside the tube.}
\end{qoptions}
\qn{8} A prism with prism angle $60^\circ$ and refractive index $\sqrt{\dfrac73}$ is given.
\qfig{p1_phy_s1_q08_1}
\begin{qoptions}
\opt{A}{Minimum possible angle of incidence so that the light ray is refracted from the second surface is $30^\circ$.}
\opt{B}{Minimum possible angle of incidence, so that the light ray is refracted from the second surface is $60^\circ$.}
\opt{C}{Maximum possible deviation produced by prism is $30^\circ$.}
\opt{D}{Maximum possible deviation produced by prism is $60^\circ$.}
\end{qoptions}
\sect{SECTION-I(ii)}
\qn{9} A particle is performing SHM $x=A\sin(\omega t+\phi)$ whose time period is $T$. List-I contains events related to the motion of particle and List-II contains the minimum time after which that event happen. Choose the correct match from option if $\phi=0^\circ$
{\small\setlength{\tabcolsep}{3pt}\renewcommand{\arraystretch}{1.25}\begin{center}\begin{tabular}{|L{0.75cm}|L{3.6cm}|L{0.7cm}|L{1.0cm}|}
\hline
\multicolumn{2}{|c|}{List--I} & \multicolumn{2}{c|}{List--II}\\
\hline
(I) & Particle comes to rest & (P) & $T/4$\\
\hline
(II) & Velocity is maximum & (Q) & $T/6$\\
\hline
(III) & Kinetic energy = Potential energy & (R) & $3T/8$\\
\hline
(IV) & Velocity is half of maximum velocity & (S) & $T/24$\\
\hline
\multicolumn{2}{|c|}{} & (T) & $T/2$\\
\hline
\multicolumn{2}{|c|}{} & (U) & $T/8$\\
\hline
\end{tabular}\end{center}}
\begin{qoptions}
\opt{A}{I $\to$ P;II $\to$ T;III $\to$ U;IV $\to$ Q}
\opt{B}{I $\to$ Q;II $\to$ P;III $\to$ R;IV $\to$ S}
\opt{C}{I $\to$ R;II $\to$ U;III $\to$ U;IV $\to$ Q}
\opt{D}{I $\to$ S;II $\to$ Q;III $\to$ R;IV $\to$ S}
\end{qoptions}
\qn{10} A particle is performing SHM $x=A\sin(\omega t+\phi)$ whose time period is $T$. List-I contains events related to the motion of particle and List-II contains the minimum time after which that event happen. Choose the correct match from option if $\phi=\pi/4$
{\small\setlength{\tabcolsep}{3pt}\renewcommand{\arraystretch}{1.25}\begin{center}\begin{tabular}{|L{0.75cm}|L{3.6cm}|L{0.7cm}|L{1.0cm}|}
\hline
\multicolumn{2}{|c|}{List--I} & \multicolumn{2}{c|}{List--II}\\
\hline
(I) & Particle comes to rest & (P) & $T/4$\\
\hline
(II) & Velocity is maximum & (Q) & $T/6$\\
\hline
(III) & Kinetic energy = Potential energy & (R) & $3T/8$\\
\hline
(IV) & Velocity is half of maximum velocity & (S) & $T/24$\\
\hline
\multicolumn{2}{|c|}{} & (T) & $T/2$\\
\hline
\multicolumn{2}{|c|}{} & (U) & $T/8$\\
\hline
\end{tabular}\end{center}}
\begin{qoptions}
\opt{A}{I $\to$ P;II $\to$ T;III $\to$ U;IV $\to$ Q}
\opt{B}{I $\to$ U;II $\to$ R;III $\to$ P;IV $\to$ S}
\opt{C}{I $\to$ R;II $\to$ Q;III $\to$ S;IV $\to$ T}
\opt{D}{I $\to$ S;II $\to$ Q;III $\to$ R;IV $\to$ T}
\end{qoptions}
\qn{11} $ABC$ is a right-angled prism kept in air. A ray (1) is incident on the face $AB$ along the normal. Refractive index of the material of prism is the minimum value that will be required so that ray (1) undergoes total internal reflection at the face $AC$. Another ray (2) is incident on the face $AB$ such that it emerges from face $AC$ along the normal to $AC$. A third ray (3) falls on the face $BC$ and emerges from face $AC$ such that its angle of emergence is the same as that of incidence. Assuming light (1), (2) and (3) have the same wavelength, then match the following.
\qfig{p1_phy_s2_q11_1}
{\small\setlength{\tabcolsep}{3pt}\renewcommand{\arraystretch}{1.25}\begin{center}\begin{tabular}{|L{0.75cm}|L{3.6cm}|L{0.7cm}|L{1.0cm}|}
\hline
\multicolumn{2}{|c|}{List-I} & \multicolumn{2}{c|}{List-II}\\
\hline
(I) & Refractive index of the material of prism is & (P) & 120\\
\hline
(II) & Angle of incidence in degree of ray (2) is & (Q) & 90\\
\hline
(III) & Deviation in degree suffered by ray (2) is & (R) & 2\\
\hline
(IV) & Deviation in degree suffered by ray (3) is & (S) & 60\\
\hline
\multicolumn{2}{|c|}{} & (T) & 1.5\\
\hline
\multicolumn{2}{|c|}{} & (U) & 4\\
\hline
\end{tabular}\end{center}}
\begin{qoptions}
\opt{A}{I $\to$ U;II $\to$ Q;III $\to$ S;IV $\to$ P}
\opt{B}{I $\to$ R;II $\to$ P;III $\to$ S;IV $\to$ P}
\opt{C}{I $\to$ U;II $\to$ S;III $\to$ P;IV $\to$ T}
\opt{D}{I $\to$ R;II $\to$ Q;III $\to$ S;IV $\to$ P}
\end{qoptions}
\qn{12} $ABC$ is a right-angled prism kept in air. A ray (1) is incident on the face $AB$ along the normal. Refractive index of the material of prism is the minimum value that will be required so that ray (1) undergoes total internal reflection at the face $AC$. Another ray (2) is incident on the face $AB$ such that it emerges from face $AC$ along the normal to $AC$. A third ray (3) falls on the face $BC$ and emerges from face $AC$ such that its angle of emergence is the same as that of incidence. Assuming light (1), (2) and (3) have the same wavelength, then match the following. [Take : $\sin^{-1}(0.29)\approx17^\circ$, $\sin15^\circ=0.258$]
\qfig{p1_phy_s2_q12_1}
{\small\setlength{\tabcolsep}{3pt}\renewcommand{\arraystretch}{1.25}\begin{center}\begin{tabular}{|L{0.75cm}|L{3.6cm}|L{0.7cm}|L{1.0cm}|}
\hline
\multicolumn{2}{|c|}{List-I} & \multicolumn{2}{c|}{List-II}\\
\hline
(I) & Refractive index of the material of prism is & (P) & 120\\
\hline
(II) & Angle of incidence in degree of ray (2) is & (Q) & 90\\
\hline
(III) & Deviation in degree suffered by ray (2) is & (R) & 2\\
\hline
(IV) & Deviation in degree suffered by ray (3) is & (S) & 30\\
\hline
\multicolumn{2}{|c|}{} & (T) & 1.15\\
\hline
\multicolumn{2}{|c|}{} & (U) & 4\\
\hline
\end{tabular}\end{center}}
\begin{qoptions}
\opt{A}{I $\to$ U;II $\to$ Q;III $\to$ S;IV $\to$ P}
\opt{B}{I $\to$ T;II $\to$ Q;III $\to$ S;IV $\to$ U}
\opt{C}{I $\to$ U;II $\to$ S;III $\to$ P;IV $\to$ T}
\opt{D}{I $\to$ R;II $\to$ Q;III $\to$ S;IV $\to$ P}
\end{qoptions}
\sect{SECTION-II}
\qn{1} A uniform chain of mass $M$ and length $L$ rests on a rough horizontal table such that one of its ends hangs over the edge of the table. Chain begins to slide if length of overhanging part is at least equal to two third of the total length of the chain. Chain is released and starts its motion from a position where two third of its total length is overhanging. If work done by friction force acting on the chain by the moment it completely slides off the table is given by $\dfrac{-\mu MgL}{N}$ then find the value of $N$. Coefficient of friction between chain and rough horizontal surface is $\mu$ and acceleration due to gravity is $g$.
\qfig{p1_phy_s3_q01_1}
\qn{2} The ends of a spring are attached to two blocks of masses 5 kg and 3 kg as shown in figure. The 5 kg block rests on a horizontal surface and the 3 kg block which is vertically above it is in equilibrium with the spring compressed by 1 cm. The 3 kg block is pushed down further by $x$ cm and released. Find the minimum value of $x$ so that the 5 kg block may be lifted off the ground during the subsequent motion.
\qfig{p1_phy_s3_q02_1}
\qn{3} A thin plano convex lens $A$ has material of refractive index $\mu_A=1.8$ and it's curved surface has radius of curvature $R$. A thin layer of transparent material $B$ is laid over the curved surface of $A$. The refractive index of $B$ is 1.2 and the curved surface of $B$ that is not touching $A$ has radius of curvature $\dfrac R2$. This surface of $B$ is silvered. A point object `$O$' is kept at a distance of 10 cm on the principal axis of the system (above the plane surface) and it's image is formed at a distance of 40 cm above the plane surface. Find $R$ (in cm).
\qfig{p1_phy_s3_q03_1}
\qn{4} A convex lens of focal length 20 cm and another plano-convex lens of focal length 40 cm are placed co-axially. The plane surface of the plano convex lens is silvered. An object $O$ is kept on the principal axis at a distance of 10 cm from the convex lens as shown in figure. Find the distance $d$ (in cm) between two lenses so that final image is formed on object itself.
\qfig{p1_phy_s3_q04_1}
\qn{5} Four small particles charged with equal positive charge $+Q$ each are arranged at the four corners of a horizontal square of side $a$. A particle having charge $+q$ and mass $m$ is placed at a height $h$ above the centre of the square. Condition under which particle at height $h$ is in equilibrium is given by
\[m=\frac{Qqh}{g\pi\varepsilon_0\left(h^2+\frac{a^2}{N_1}\right)^{N_2}}.\]
Fill the value of $(N_1+N_2)$. (Acceleration due to gravity is `$g$')
\qfig{p1_phy_s3_q05_1}
\qn{6} A particle performs SHM along $x$-axis according to the relation $x=3\sin(100t)+8\cos^2(50t)$ where `$x$' is in cm and `$t$' is in second. Find the maximum possible distance (in cm) of the particle from origin during its motion.
\raggedcolumns
\partsec{2}{CHEMISTRY}
\sect{SECTION-I(i)}
\qn{1} Which of the following is/are correctly match with their equilibrium constant ($K_{\mathrm{eq}}$)
\begin{qoptions}
\opt{A}{\qfigi{p1_chem_s1_q01_A}}
\opt{B}{\qfigi{p1_chem_s1_q01_B}}
\opt{C}{\qfigi{p1_chem_s1_q01_C}}
\opt{D}{\qfigi{p1_chem_s1_q01_D}}
\end{qoptions}
\qn{2} Which of the following reaction result the formation of a pair of diastereomers?
\begin{qoptions}
\opt{A}{\qfigi{p1_chem_s1_q02_A}}
\opt{B}{\qfigi{p1_chem_s1_q02_B}}
\opt{C}{\qfigi{p1_chem_s1_q02_C}}
\opt{D}{\qfigi{p1_chem_s1_q02_D}}
\end{qoptions}
\qn{3} \qfigi{p1_chem_s1_q03_1} $A+B$\par
Identify relation between $A$ and $B$. If both $A$ and $B$ give Bromine water test
\begin{qoptions}
\opt{A}{Stereoisomers}
\opt{B}{Structural isomers}
\opt{C}{Enantiomers}
\opt{D}{Diastereomers}
\end{qoptions}
\qn{4} Choose the correct statement among the following :
\begin{qoptions}
\opt{A}{Primitive cubic unit cell involves two dimensional hexagonal closed packed array of atoms}
\opt{B}{Four hexagonal closed packed two dimensional layers of atoms are involved in one FCC unit cell}
\opt{C}{Three hexagonal close packed two dimensional layers of atoms are involved in one HCP unit cell}
\opt{D}{Four square close packed two dimensional layers of atoms are involved in one BCC unit cell}
\end{qoptions}
\qn{5} Hydrocarbon contains $\dfrac{600}{7}$\% by mass carbon and its one molecule weighs $93.33\times10^{-24}$ g. Which of the following statement(s) is/are correct\ $[N_A=6.0\times10^{23}]$
\begin{qoptions}
\opt{A}{Empirical formula is $\mathrm{CH}$}
\opt{B}{Empirical formula is $\mathrm{CH_2}$}
\opt{C}{Molecular formula of $\mathrm{C_3H_8}$}
\opt{D}{Molecular formula is $\mathrm{C_4H_8}$}
\end{qoptions}
\qn{6} Which of the following given order is \textbf{CORRECT} with respect to bond angle?
\begin{qoptions}
\opt{A}{$\mathrm{Cl{-}\widehat{Se}{-}Cl}>\mathrm{Cl{-}\widehat{Se}{-}O}$ (In $\mathrm{SeOCl_2}$)}
\opt{B}{$\mathrm{F_{eq}{-}\widehat{S}{-}F_{eq}}$(in $\mathrm{SF_4}$)$>\mathrm{F_{eq}{-}\widehat{Sb}{-}F_{eq}}$ (In $\mathrm{SbF_4^{-}}$)}
\opt{C}{$\mathrm{F_{eq}{-}\widehat{S}{-}F_{eq}}$(in $\mathrm{SF_4}$)$<\mathrm{F_{eq}{-}\widehat{Sb}{-}F_{eq}}$ (In $\mathrm{SbF_4^{-}}$)}
\opt{D}{$\mathrm{Cl{-}\widehat{Se}{-}Cl}<\mathrm{Cl{-}\widehat{Se}{-}O}$ (In $\mathrm{SeOCl_2}$)}
\end{qoptions}
\qn{7} Which of the following species is/are planar and polar?
\begin{qoptions}
\opt{A}{$\mathrm{NO_2}$}
\opt{B}{$\mathrm{ClO_2}$}
\opt{C}{$\mathrm{XeF_5^{-}}$}
\opt{D}{$\mathrm{ClF_3}$}
\end{qoptions}
\qn{8} Choose the INCORRECT order of the indicated property given in the options.
{\small\setlength{\tabcolsep}{3pt}\renewcommand{\arraystretch}{1.25}\begin{center}\begin{tabular}{|L{0.6cm}|L{3.3cm}|c|L{2.3cm}|}
\hline
(A) & $\mathrm{He<Ne<Kr<Ar}$ & : & Clathrate compound formation tendency order\\
\hline
(B) & $\mathrm{NaO_2}<\mathrm{KO_2}<\mathrm{RbO_2}<\mathrm{CsO_2}$ & : & Thermal stability order\\
\hline
(C) & $\mathrm{BeF_2<BaF_2}$ & : & Solubility order\\
\hline
(D) & $\mathrm{LiF>LiCl>LiI}$ & : & Melting point order\\
\hline
\end{tabular}\end{center}}
\sect{SECTION-I(ii)}
\qn{9} \textbf{Answer the following by approximately matching the lists based on the information given in the paragraph.}\par
For an $n^{\mathrm{th}}$ order reaction $\mathrm A\to$ products,\par
The rate law is $\dfrac{-dC_A}{dt}=kC_A^{\,n}$. Where $C_A$ = concentration of A at time `$t$'. Let $C_{A_0}$ is initial concentration of A and and $\alpha$ = degree of dissociation of reactant `A'
{\small\setlength{\tabcolsep}{3pt}\renewcommand{\arraystretch}{1.25}\begin{center}\begin{tabular}{|L{0.85cm}|L{1.2cm}|L{0.65cm}|L{3.5cm}|}
\hline
 & List-I & & List-II\\
\hline
(I) & $n=0$ & (P) & $\alpha=\dfrac{C_{A_0}-C_A}{C_{A_0}}$\\
\hline
(II) & $n=\dfrac12$ & (Q) & $\alpha=1-e^{-kt}$\\
\hline
(III) & $n=1$ & (R) & $\alpha=\dfrac{kt}{C_{A_0}}$\\
\hline
(IV) & $n=2$ & (S) & $\alpha=\dfrac{C_{A_0}kt}{(1+C_{A_0}kt)}$\\
\hline
 & & (T) & $\alpha$ = independent of initial concentration of reactant\\
\hline
 & & (U) & Graph between `$\alpha$' and `$t$' is a straight line passing through origin\\
\hline
\end{tabular}\end{center}}
Which of the following options has the correct combination considering List-I and List-II ?
\begin{qoptions}
\opt{A}{II $\to$ P, R}
\opt{B}{I $\to$ P,R,T}
\opt{C}{III $\to$ P,Q,U}
\opt{D}{III $\to$ P,Q,T}
\end{qoptions}
\qn{10} \textbf{Answer the following by approximately matching the lists based on the information given in the paragraph.}\par
For an $n^{\mathrm{th}}$ order reaction $\mathrm A\to$ products,\par
The rate law is $\dfrac{-dC_A}{dt}=kC_A^{\,n}$. Where $C_A$ = concentration of A at time `$t$'. Let $C_{A_0}$ is initial concentration of A and and $\alpha$ = degree of dissociation of reactant `A'
{\small\setlength{\tabcolsep}{3pt}\renewcommand{\arraystretch}{1.25}\begin{center}\begin{tabular}{|L{0.85cm}|L{1.2cm}|L{0.65cm}|L{3.5cm}|}
\hline
 & List-I & & List-II\\
\hline
(I) & $n=0$ & (P) & $\alpha=\dfrac{C_{A_0}-C_A}{C_{A_0}}$\\
\hline
(II) & $n=\dfrac12$ & (Q) & $\alpha=1-e^{-kt}$\\
\hline
(III) & $n=1$ & (R) & $\alpha=\dfrac{kt}{C_{A_0}}$\\
\hline
(IV) & $n=2$ & (S) & $\alpha=\dfrac{C_{A_0}kt}{(1+C_{A_0}kt)}$\\
\hline
 & & (T) & $\alpha$ = independent of initial concentration of reactant\\
\hline
 & & (U) & Graph between `$\alpha$' and `$t$' is a straight line passing through origin\\
\hline
\end{tabular}\end{center}}
Which of the following options has the correct combination considering List-I and List-II ?
\begin{qoptions}
\opt{A}{I $\to$ P,Q}
\opt{B}{IV $\to$ S,T}
\opt{C}{IV $\to$ P,S}
\opt{D}{III $\to$ Q,U}
\end{qoptions}
\qn{11} {\small\setlength{\tabcolsep}{3pt}\renewcommand{\arraystretch}{1.25}\begin{center}\begin{tabular}{|L{0.8cm}|L{3.0cm}|L{0.65cm}|L{1.6cm}|}
\hline
\multicolumn{2}{|c|}{\parbox{3.3cm}{\centering List-I\\Dimerization/ polymerization by}} & \multicolumn{2}{c|}{List-II}\\
\hline
(I) & $3c-2e^-$ & (P) & $\mathrm{BeX_2}$\\
\hline
(II) & Co-ordinate bond & (Q) & $\mathrm{AlX_3}$\\
\hline
(III) & Odd $e^-$ & (R) & $\mathrm{CX_3}$\\
\hline
(IV) & No dimerization & (S) & $\mathrm{BX_3}$\\
\hline
\multicolumn{2}{|c|}{} & (T) & $\mathrm{IX_3}$\\
\hline
\multicolumn{2}{|c|}{} & (U) & $\mathrm{ClO_3}$\\
\hline
\end{tabular}\end{center}}
Select correct match when $X=\mathrm F$.
\begin{qoptions}
\opt{A}{II $\to$ R}
\opt{B}{I $\to$ Q}
\opt{C}{III $\to$ U}
\opt{D}{IV $\to$ S}
\end{qoptions}
\qn{12} {\small\setlength{\tabcolsep}{3pt}\renewcommand{\arraystretch}{1.25}\begin{center}\begin{tabular}{|L{0.8cm}|L{3.0cm}|L{0.65cm}|L{1.6cm}|}
\hline
\multicolumn{2}{|c|}{\parbox{3.3cm}{\centering List-I\\Dimerization/ polymerization by}} & \multicolumn{2}{c|}{List-II}\\
\hline
(I) & $3c-2e^-$ & (P) & $\mathrm{BeX_2}$\\
\hline
(II) & Co-ordinate bond & (Q) & $\mathrm{AlX_3}$\\
\hline
(III) & Odd $e^-$ & (R) & $\mathrm{CX_3}$\\
\hline
(IV) & No dimerization & (S) & $\mathrm{BX_3}$\\
\hline
\multicolumn{2}{|c|}{} & (T) & $\mathrm{IX_3}$\\
\hline
\multicolumn{2}{|c|}{} & (U) & $\mathrm{ClO_3}$\\
\hline
\end{tabular}\end{center}}
Select correct match when $X=\mathrm{Cl}$
\begin{qoptions}
\opt{A}{I $\to$ P}
\opt{B}{III $\to$ R}
\opt{C}{II $\to$ Q}
\opt{D}{IV $\to$ T}
\end{qoptions}
\qfullbegin
\sect{SECTION-II}
\qn{1} Number of following compounds which can't be synthesised in good yield by Williamson's reaction
\qfig{p1_chem_s3_q01_1}
\qfullend
\qn{2} Number of compounds which are more reactive than $\mathrm{CH_3{-}CH(CH_3){-}Cl}$ by $\mathrm{S_N^{\,2}}$ mechanism?
\qfig{p1_chem_s3_q02_2}
\qn{3} How many of the following substrate can give E2 elimination reaction on treatment with base / $\Delta$
\qfig{p1_chem_s3_q03_1}
\qn{4} What should be the minimum pressure (in atm) applied on 0.1 M glucose solution, filled in semi-permiable membrane, so as to filter away 750 mL water, leaving 250 mL of glucose solution at 300 K.\\\relax
[Given : $R=0.08$ atm L$\mathrm{K^{-1}}$ mol$^{-1}$]
\qn{5} An aq. solution of 0.5M $\mathrm{KMnO_4}$ is divided into two parts. One part of it requires 125 ml of 1.5M aq. solution of oxalate ions in acidic medium, while another part requires 270 ml of 0.5M aq. solution of iodide ions in neutral medium which are converted into $\mathrm{I_2}$ only. Calcualte total volume (mL) of the initial $\mathrm{KMnO_4}$ solution.
\qn{6} If s and p-orbital mixing (s-p mixing) is operative for all species given, then how many following species is/are paramagnetic in nature ?\\
$\mathrm{O_2^-}$, $\mathrm{N_2}$, $\mathrm{C_2}$, $\mathrm{CN^-}$, $\mathrm{B_2}$, $\mathrm{O_2^{2+}}$
\raggedcolumns
\partsec{3}{MATHEMATICS}
\sect{SECTION-I(i)}
\qn{1} Consider $f:R\to R$ satisfying the equation $f(x+y)=f(x-y)+f(x).f(y)$ $\forall\ x,y\in R$.\
Identify the correct statement(s)
\begin{qoptions}
\opt{A}{$f(2014)+f(2015)=f(2016)$}
\opt{B}{$\sin^{-1}(f(2))+\cos^{-1}(f(3))=\dfrac\pi2$}
\opt{C}{$f(2014)-f(2015)=f(2016)$}
\opt{D}{$\sin^{-1}(f(3))+\cos^{-1}(f(2))=\dfrac\pi2$}
\end{qoptions}
\qn{2} $ABC$ is a triangle in which $\angle C=2\angle B$. $D$ is a point on $BC$ such that $AD$ bisects $\angle BAC$ and $AC=BD$. Then which of the following is/are true
\begin{qoptions}
\opt{A}{$\cos B+\cos C=\dfrac{\sqrt5}{2}$}
\opt{B}{$\cos B-\cos C=\dfrac12$}
\opt{C}{$\dfrac{\sin A}{\sin C}=1$}
\opt{D}{$\dfrac{\sin A}{\sin C}=\dfrac{\sqrt5+1}{4}$}
\end{qoptions}
\qn{3} For the figure given below if $\tan\alpha=\dfrac mn$, where $m$ and $n$ are relatively prime positive integers then
\qfig{p1_math_s1_q03_1}
\begin{qoptions}
\opt{A}{$m+n=463$}
\opt{B}{$n-m=122$}
\opt{C}{$2m-n=41$}
\opt{D}{$m+n=465$}
\end{qoptions}
\qn{4} The value of $\displaystyle\lim_{t\to1^-}(1-t)\sum_{r=1}^{\infty}\dfrac{t^r}{1+t^r}$ is
\begin{qoptions}
\opt{A}{$\ln2$}
\opt{B}{$\ln3$}
\opt{C}{$\ln\left(\dfrac{2e}{1+e}\right)$}
\opt{D}{$\dfrac12$}
\end{qoptions}
\qn{5} Which of the following are even functions in their domain
\begin{qoptions}
\opt{A}{$y(x)=\dfrac{(2\sin x-1)(\sin^2x-1)}{(2\sin x-1)}$}
\opt{B}{$y(x)=f\mathrm{ogoh}(x)$ (where $f$ is neither even nor odd, $g$ is even and $h$ is an odd function, all functions are defined for all $x\in R$)}
\opt{C}{$y(x)$ is defined for all real $x$ and $y(x)=|x-1|$ for $x\in[0,2]$, and $y(x)$ is periodic with period 2.}
\opt{D}{$y(x)=\dfrac{1+e^{x-1}\left(\cos(1-x)+\cos(1+x)\right)+e^{2x}}{e^x}$}
\end{qoptions}
\qn{6} For the equation $\cos^{-1}\left(\dfrac{x^2-1}{x^2+1}\right)+\sin^{-1}\left(\dfrac{2x}{x^2+1}\right)+\tan^{-1}\left(\dfrac{2x}{x^2-1}\right)=\dfrac{2\pi}3$, which of the following is/are correct.
\begin{qoptions}
\opt{A}{The number of solution is 2}
\opt{B}{It has at least one irrational solution}
\opt{C}{The sum of solution is positive}
\opt{D}{The sum of solution is negative}
\end{qoptions}
\qn{7} In the diagram shown, $AB$ and $AC$ are equal sides of an isosceles triangle $ABC$. The equilateral triangle $DEF$ is inscribed. If $\angle ADF$ be $\alpha$, $\angle BED$ be $\beta$ and $\angle CFE$ be $\gamma$, then \textbf{CORRECT} relation is/are
\qfig{p1_math_s1_q07_1}
\begin{qoptions}
\opt{A}{$\beta=\dfrac{\alpha+\gamma}2$}
\opt{B}{$\alpha=\dfrac{\beta+\gamma}2$}
\opt{C}{$\beta=\dfrac{\alpha-\gamma}2$}
\opt{D}{$\alpha=\dfrac{\beta-\gamma}2$}
\end{qoptions}
\qn{8} If $\alpha$ ($\alpha>0$) is a root of $x^{11}-x^{10}+x^8-x^7+x^5-x^4+x^2-x-20=0$ then $\alpha^{12}$ is greater than
\begin{qoptions}
\opt{A}{59}
\opt{B}{61}
\opt{C}{58}
\opt{D}{57}
\end{qoptions}
\sect{SECTION-I(ii)}
\qn{9} Match List-I with List-II and select the correct answer using the code given below the list.
{\small\setlength{\tabcolsep}{3pt}\renewcommand{\arraystretch}{1.25}\begin{center}\begin{tabular}{|L{0.8cm}|L{4.2cm}|L{0.65cm}|L{0.75cm}|}
\hline
\multicolumn{2}{|c|}{\parbox{4.6cm}{\centering List-I\\Equations}} & \multicolumn{2}{c|}{\parbox{1.5cm}{\centering\vphantom{\rule{0pt}{2ex}}List-II\\No. of solutions\vphantom{\rule[-1ex]{0pt}{1ex}}}}\\
\hline
(I) & $x^2\tan x=1$\ \ $x\in\left(\dfrac\pi2,\dfrac{3\pi}2\right)$ & (P) & 5\\
\hline
(II) & $2^{\cos x}=|\sin x|$, $x\in[0,2\pi]$ & (Q) & 1\\
\hline
(III) & If $f(x)$ is a polynomial of degree 5 with real coefficients such that $f(|x|)=0$ has 8 real roots then, roots of $f(x)=0$ & (R) & 3\\
\hline
(IV) & $7^{|x|}\left(|5-|x||\right)=1$ & (S) & 4\\
\hline
\multicolumn{2}{|c|}{} & (T) & 0\\
\hline
\multicolumn{2}{|c|}{} & (U) & 10\\
\hline
\end{tabular}\end{center}}
Which of the following is the only \textbf{CORRECT} combination ?
\begin{qoptions}
\opt{A}{I $\to$ P}
\opt{B}{II $\to$ Q}
\opt{C}{III $\to$ R}
\opt{D}{IV $\to$ S}
\end{qoptions}
\qn{10} Match List-I with List-II and select the correct answer using the code given below the list.
{\small\setlength{\tabcolsep}{3pt}\renewcommand{\arraystretch}{1.25}\begin{center}\begin{tabular}{|L{0.8cm}|L{4.2cm}|L{0.65cm}|L{0.75cm}|}
\hline
\multicolumn{2}{|c|}{\parbox{4.6cm}{\centering List-I\\Equations}} & \multicolumn{2}{c|}{\parbox{1.5cm}{\centering\vphantom{\rule{0pt}{2ex}}List-II\\No. of solutions\vphantom{\rule[-1ex]{0pt}{1ex}}}}\\
\hline
(I) & $x^2\tan x=1$\ \ $x\in\left(\dfrac\pi2,\dfrac{3\pi}2\right)$ & (P) & 5\\
\hline
(II) & $2^{\cos x}=|\sin x|$, $x\in[0,2\pi]$ & (Q) & 1\\
\hline
(III) & If $f(x)$ is a polynomial of degree 5 with real coefficients such that $f(|x|)=0$ has 8 real roots then, roots of $f(x)=0$ & (R) & 3\\
\hline
(IV) & $7^{|x|}\left(|5-|x||\right)=1$ & (S) & 4\\
\hline
\multicolumn{2}{|c|}{} & (T) & 0\\
\hline
\multicolumn{2}{|c|}{} & (U) & 10\\
\hline
\end{tabular}\end{center}}
Which of the following is the only \textbf{INCORRECT} combination ?
\begin{qoptions}
\opt{A}{I $\to$ R}
\opt{B}{II $\to$ S}
\opt{C}{III $\to$ P}
\opt{D}{IV $\to$ Q}
\end{qoptions}
\qn{11} Match List-I with List-II and select the correct answer using the code given below the list.
{\small\setlength{\tabcolsep}{3pt}\renewcommand{\arraystretch}{1.25}\begin{center}\begin{tabular}{|L{0.8cm}|L{4.6cm}|L{0.65cm}|L{0.4cm}|}
\hline
\multicolumn{2}{|c|}{List-I} & \multicolumn{2}{c|}{List-II}\\
\hline
(I) & Number of solutions of $4(\sin^6\theta+\cos^6\theta)=10-18\sin\theta\cos\theta$ in $[0,3\pi]$, is & (P) & 1\\
\hline
(II) & Number of possible pairs $(x,y)$ satisfying $x^2-6x+10-\sin^2y=0$, where $y\in[0,2\pi]$ and $x\in R$, is & (Q) & 2\\
\hline
(III) & Number of integral values of $x$ satisfying $\sin x<\dfrac12$ in $[0,2\pi]$, is & (R) & 3\\
\hline
(IV) & Number of solutions satisfying $\tan x+\sec x=2\cos x$ in $[0,4\pi]$, is & (S) & 4\\
\hline
\multicolumn{2}{|c|}{} & (T) & 5\\
\hline
\multicolumn{2}{|c|}{} & (U) & 6\\
\hline
\end{tabular}\end{center}}
Which of the following is the only \textbf{CORRECT} combination ?
\begin{qoptions}
\opt{A}{I $\to$ Q}
\opt{B}{I $\to$ R}
\opt{C}{II $\to$ P}
\opt{D}{II $\to$ R}
\end{qoptions}
\qn{12} Match List-I with List-II and select the correct answer using the code given below the list.
{\small\setlength{\tabcolsep}{3pt}\renewcommand{\arraystretch}{1.25}\begin{center}\begin{tabular}{|L{0.8cm}|L{4.6cm}|L{0.65cm}|L{0.4cm}|}
\hline
\multicolumn{2}{|c|}{List-I} & \multicolumn{2}{c|}{List-II}\\
\hline
(I) & Number of solutions of $4(\sin^6\theta+\cos^6\theta)=10-18\sin\theta\cos\theta$ in $[0,3\pi]$, is & (P) & 1\\
\hline
(II) & Number of possible pairs $(x,y)$ satisfying $x^2-6x+10-\sin^2y=0$, where $y\in[0,2\pi]$ and $x\in R$, is & (Q) & 2\\
\hline
(III) & Number of integral values of $x$ satisfying $\sin x<\dfrac12$ in $[0,2\pi]$, is & (R) & 3\\
\hline
(IV) & Number of solutions satisfying $\tan x+\sec x=2\cos x$ in $[0,4\pi]$, is & (S) & 4\\
\hline
\multicolumn{2}{|c|}{} & (T) & 5\\
\hline
\multicolumn{2}{|c|}{} & (U) & 6\\
\hline
\end{tabular}\end{center}}
Which of the following is the only \textbf{CORRECT} combination ?
\begin{qoptions}
\opt{A}{III $\to$ S}
\opt{B}{III $\to$ U}
\opt{C}{IV $\to$ T}
\opt{D}{IV $\to$ S}
\end{qoptions}
\sect{SECTION-II}
\qn{1} Let $S_n=\displaystyle\sum_{k=0}^{n}\tan^{-1}\dfrac{49}{49k^2+63k+57}$ then $\cot\left(\displaystyle\lim_{n\to\infty}S_n\right)$ is equal to
\qn{2} In a triangle $ABC$, medians $AD$ and $CE$ are drawn. If $AD=5$, $\angle DAC=\dfrac\pi8$ and $\angle ACE=\dfrac\pi4$ then the area of the triangle $ABC$ is equal to
\qn{3} If $\displaystyle\lim_{x\to0^+}\left(\dfrac{\sin^2x}{\left[\frac x\pi\right]+\frac{x^2}{\pi^2}}\right)\times\dfrac{\sin(\sin x)-\sin x}{ax^5+bx^3+c}=\dfrac{-\pi^2}{12}$. Then value of $b+c$ is\
(Where $[.]$ denotes greatest integer function)
\qn{4} If the value of $\displaystyle\lim_{x\to\infty}x^2\tan\left(\ln\left(\sqrt{\cos\dfrac\pi x}\right)\right)$ is equal to $-k\pi^2$, then $k$ is equal to
\qn{5} Let \[\begin{aligned}f(\theta)=\dfrac1{\tan^9\theta}\Big(&(1+\tan\theta)^{10}+(2+\tan\theta)^{10}\\&+\ldots.(20+\tan\theta)^{10}\Big)-20\tan\theta.\end{aligned}\] The value of $\displaystyle\lim_{\theta\to\frac\pi2^-}f(\theta)$ is
\qn{6} The number of solution of equation $2\sin2x+\sin x(1-2\cos2x)-4=0$ in the interval $(0,\pi)$ :
\raggedcolumns
\newpage
\papersec{2}
\partsec{1}{PHYSICS}
\sect{SECTION-I}
\qn{1} A ball of radius $R$ carries a positive charge whose volume charge density depends only on separation $r$ from the ball's centre as $\rho=\rho_0\left(1-\dfrac rR\right)$ where $\rho_0$ is constant. The electric field at distance $r_1$ from centre (For $r_1<R$) is :-
\begin{qoptions}
\opt{A}{$\dfrac{\rho_0r_1}{3\epsilon_0}\left(1-\dfrac{3r_1}{4R}\right)$}
\opt{B}{$\dfrac{\rho_0r_1}{3\epsilon_0}\left(1-\dfrac{r_1}{R}\right)$}
\opt{C}{$\dfrac{\rho_0r_1}{2\epsilon_0}\left(1-\dfrac{r_1}{2R}\right)$}
\opt{D}{None of these}
\end{qoptions}
\qn{2} A particle of mass $m$ and carrying charge $-q_1$ is moving around a fixed charge $+q_2$ along a circular path of radius $r$. Find period of revolution of the charge $-q_1$ about $+q_2$.
\begin{qoptions}
\opt{A}{$\sqrt{\dfrac{16\pi^3\epsilon_0mr^3}{q_1q_2}}$}
\opt{B}{$\sqrt{\dfrac{8\pi^3\epsilon_0mr^3}{q_1q_2}}$}
\opt{C}{$\sqrt{\dfrac{q_1q_2}{16\pi^3\epsilon_0mr^3}}$}
\opt{D}{Zero}
\end{qoptions}
\qn{3} In a region, electric field depends on $X$-axis as $E=E_0x^2$. There is a cube of edge $a$ as shown. Then find the charge enclosed in that cube.
\qfig{p2_phy_s1_q03_1}
\begin{qoptions}
\opt{A}{$5\epsilon_0a^4E_0$}
\opt{B}{$3\epsilon_0a^4E_0$}
\opt{C}{$4\epsilon_0a^4E_0$}
\opt{D}{Zero}
\end{qoptions}
\qn{4} A particle is performing SHM and its equation is $x=10\sin\left(4t-\dfrac\pi3\right)$, where time ($t$) is in second and $x$ is the position from the mean position (in metre). The speed of particle at the instant when its displacement from mean position is 6m, is:-
\begin{qoptions}
\opt{A}{$\dfrac{16\pi}3$ m/s}
\opt{B}{$\dfrac{10\pi}3$ m/s}
\opt{C}{18 m/s}
\opt{D}{32 m/s}
\end{qoptions}
\qn{5} A spherical cavity is made in a uniformly charged non-conducting sphere and a positive charge is released from position $P$ as shown in figure. Then path followed by particle is :-
\qfig{p2_phy_s1_q05_1}
\begin{qoptions}
\opt{A}{A straight line along radially inward of cavity}
\opt{B}{A straight line along radially outward of cavity}
\opt{C}{A straight line along positive $X$-direction}
\opt{D}{A straight line along radially outward of sphere}
\end{qoptions}
\qn{6} A linear object is placed along the axis of a mirror as shown in the figure. If $f$ is the focal length of the mirror then the length of image is
\qfig{p2_phy_s1_q06_1}
\begin{qoptions}
\opt{A}{$\dfrac{2f}3$}
\opt{B}{$f$}
\opt{C}{$\dfrac f3$}
\opt{D}{$\dfrac{4f}3$}
\end{qoptions}
\qn{7} A pump of 200 W is lifting 2kg water per second from an average depth of 10m. The velocity with which water comes out of pump is :- (Take : $g=9.8$ m/s$^2$)
\begin{qoptions}
\opt{A}{2 m/s}
\opt{B}{2.5 m/s}
\opt{C}{4 m/s}
\opt{D}{3 m/s}
\end{qoptions}
\qn{8} A system is released from rest as shown in figure. Kinetic energy of mass $m$ when its moves distance $x$ is downward direction is. (initially both spring are unstretched and all spring are massless)
\qfig{p2_phy_s1_q08_1}
\begin{qoptions}
\opt{A}{$10\,mgx$}
\opt{B}{$\dfrac{10mgx-kx^2}{10}$}
\opt{C}{$mgx-kx^2$}
\opt{D}{$mgx-\dfrac{kx^2}2$}
\end{qoptions}
\qn{9} Two bodies of mass $m$ and $4m$ are attached by a string shown in the figure. The body of mass $m$ hanging from a string of length $\ell$ is executing simple harmonic motion with small amplitude $A$ while other body is at rest on the surface. The minimum coefficient of friction between the mass $4m$ and the horizontal surface must be :
\qfig{p2_phy_s1_q09_1}
\begin{qoptions}
\opt{A}{$\dfrac14\left(1-\dfrac{A^2}{\ell^2}\right)$}
\opt{B}{$\dfrac14\left(1+\dfrac{A^2}{\ell^2}\right)$}
\opt{C}{$\dfrac14\dfrac A\ell\cos\theta$}
\opt{D}{$\dfrac14$}
\end{qoptions}
\qn{10} A particle is displaced form $x=-6$m to $x=+6$m. A force $F$ acting on the particle during its motion is shown in figure. Graph between work done by this force ($W$) and displacement ($x$) should be :-
\qfig{p2_phy_s1_q10_1}
\begin{qoptions}
\opt{A}{\qfigi{p2_phy_s1_q10_A}}
\opt{B}{\qfigi{p2_phy_s1_q10_B}}
\opt{C}{\qfigi{p2_phy_s1_q10_C}}
\opt{D}{\qfigi{p2_phy_s1_q10_D}}
\end{qoptions}
\qn{11} A bob of mass $m$ is attached to the lower end of a rigid light rod of length $L$ and other end is fixed to a point $O$ so that the system may rotate in vertical circle about $O$ as shown in the figure. Initially the bob is at point $A$ and rod makes an angle $45^\circ$ from vertical. What minimum speed ($v$) should be imparted to the bob so that it may complete the circular track?
\qfig{p2_phy_s1_q11_1}
\begin{qoptions}
\opt{A}{$u=\sqrt{gL(2-\sqrt2)}$}
\opt{B}{$u=\sqrt{gL(2+\sqrt2)}$}
\opt{C}{$\sqrt{gL(3+\sqrt2)}$}
\opt{D}{$\sqrt{gL(3-\sqrt2)}$}
\end{qoptions}
\qn{12} A force acting on a particle moving in the $xy$-plane is given by $\vec F=\left(2y\hat i+x^2\hat j\right)$ N, where $x$ and $y$ are in m. The particle moves along a straight line from the origin to $(5,5)$. The work done by $F$ is :-
\begin{qoptions}
\opt{A}{125 J}
\opt{B}{66.7 J}
\opt{C}{35 J}
\opt{D}{25 J}
\end{qoptions}
\qn{13} If an astronomical telescope has objective and eye-piece of focal lengths 100 cm and 4 cm respectively, then the magnifying power of the telescope for the normal vision is :
\begin{qoptions}
\opt{A}{42}
\opt{B}{25}
\opt{C}{58}
\opt{D}{104}
\end{qoptions}
\qn{14} A particle in a certain conservative force field has a potential energy given by $U=\dfrac{20yz}{x}$. The force exerted on it is :-
\begin{qoptions}
\opt{A}{$\left(\dfrac{20yz}{x^2}\right)\hat i+\left(\dfrac{20z}x\right)\hat j+\left(\dfrac{20y}x\right)\hat k$}
\opt{B}{$\left(\dfrac{20yz}{x^2}\right)\hat i-\left(\dfrac{20z}x\right)\hat j-\left(\dfrac{20y}x\right)\hat k$}
\opt{C}{$-\left(\dfrac{20yz}{x^2}\right)\hat i-\left(\dfrac{20z}x\right)\hat j-\left(\dfrac{20y}x\right)\hat k$}
\opt{D}{$-\left(\dfrac{20yz}{x^2}\right)\hat i+\left(\dfrac{20z}x\right)\hat j+\left(\dfrac{20y}x\right)\hat k$}
\end{qoptions}
\qn{15} A convex lens $A$ of focal length 20 and a concave lens $B$ of focal length 5 cm are kept along the same axis with the distance $d$ between them. If a parallel beam of light falling on $A$ leaves $B$ as a parallel beam, then distance $d$ in cm will be :-
\begin{qoptions}
\opt{A}{25}
\opt{B}{15}
\opt{C}{30}
\opt{D}{50}
\end{qoptions}
\qn{16} A positive point charge `$Q$' is located at origin and a negative point charge `$q$' is located on $x$-axis at a distance `$r$' from `$Q$'. Magnitude of charge of `$Q$' is greater than that of `$q$'. Then, which of the following may represent the $E$-$x$ graph?
\begin{qoptions}
\opt{A}{\qfigi{p2_phy_s1_q16_A}}
\opt{B}{\qfigi{p2_phy_s1_q16_B}}
\opt{C}{\qfigi{p2_phy_s1_q16_C}}
\opt{D}{\qfigi{p2_phy_s1_q16_D}}
\end{qoptions}
\qn{17} Two identical prisms with slightly different indices are located as shown. Angle $\theta$ is small. When a laser beam strikes one of the prisms perpendicular to the surface, the refracted ray is deviated by a small angle $\phi$. The difference between the indices of refraction of the prism in terms of $\theta$ and $\phi$ is :-
\qfig{p2_phy_s1_q17_1}
\begin{qoptions}
\opt{A}{$\dfrac\phi\theta$}
\opt{B}{$\dfrac\theta\phi$}
\opt{C}{$\dfrac{2\theta}\phi$}
\opt{D}{$\dfrac{2\phi}\theta$}
\end{qoptions}
\qn{18} Figure shows a graph of angle of deviation $\delta$ versus angle of incidence $i$ for a light ray striking a prism. Angle of prism is
\qfig{p2_phy_s1_q18_1}
\begin{qoptions}
\opt{A}{$30^\circ$}
\opt{B}{$45^\circ$}
\opt{C}{$60^\circ$}
\opt{D}{$75^\circ$}
\end{qoptions}
\qn{19} A ray of light is incident on a glass sphere of refractive index $\sqrt3$ at an angle of incidence $60^\circ$ as shown in figure. The total deviation suffered by the ray is
\qfig{p2_phy_s1_q19_1}
\begin{qoptions}
\opt{A}{$30^\circ$}
\opt{B}{$60^\circ$}
\opt{C}{$90^\circ$}
\opt{D}{$120^\circ$}
\end{qoptions}
\qn{20} In an astronomical telescope, objective is of focal length 20 cm. What should be tube length for magnifying power of 10, under normal magnification ?
\begin{qoptions}
\opt{A}{10 cm}
\opt{B}{22 cm}
\opt{C}{24 cm}
\opt{D}{30 cm}
\end{qoptions}
\sect{SECTION-II}
\qn{1} A sphere of radius $R$ have volume charge density give as $\rho(r)=Kr$ for $r\le R$ and 0 for $r>R$. If electric field at distance $\dfrac R2$ from centre is $\dfrac{KR^2}{N\epsilon_0}$, then $N$ will be :
\qn{2} A point object is placed at distance of 20 cm from a lens as shown in figure. The distance of final image from the lens (in cm) is \underline{\hspace{1cm}}.
\qfig{p2_phy_s2_q02_1}
\qn{3} One end of an ideal spring is fixed to a wall at origin $O$ and axis of the spring is parallel to $x$-axis. A block of mass 1 kg is attached to the free end of the spring and is performing SHM. Equation of position of the block in co-ordinate system shown in figure is $x=10+3\sin(10t)$, where $t$ is in second and $x$ in cm. Another identical block $B$, moving towards the origin with velocity 0.6 m/s collides elastically with the block $A$ at $t=0$. The new amplitude (in cm) of oscillations will be :
\qfig{p2_phy_s2_q03_1}
\qn{4} Two identical springs are attached to a small block $P$. The other ends of the springs are fixed at $A$ and $B$ such that $AB$ is vertical. When $P$ is in equilibrium the extension of top spring is 20 cm and extension of bottom spring is 10 cm. If the period of small vertical oscillations of $P$ about its equilibrium position is $\dfrac{\pi}{N\sqrt2}$ sec, then fill the value of $N$.
\qfig{p2_phy_s2_q04_1}
\qn{5} A quarter cylinder of glass of radius $R$ and refractive index 1.5 is placed on a table. A point object $P$ is kept at a distance of $mR$ from it. Find the value of $3m$ for which a ray from $P$ will emerge parallel to the table as shown in the figure.
\qfig{p2_phy_s2_q05_1}
\qn{6} The convex surface of a thin concavo-convex lens of glass of refractive index 1.5 has a radius of curvature 20 cm. The concave surface has a radius of curvature 60 cm. The convex side is silvered and placed on a horizontal surface. At what distance (in cm) should a pin be placed on optical axis so that image is formed at same place.
\qfig{p2_phy_s2_q06_1}
\qn{7} Two particles are performing simple harmonic motion with same amplitude and same frequency. When they are at same distance from mean position on opposite sides, their speeds are one fourth of maximum speed. Also at these positions their direction of velocities are opposite. If the phase difference between the two simple harmonic motions is $\dfrac{3\pi}{x}$, then find the value of $x$.
\qn{8} The distance between the object and its real image from the convex lens is 60 cm and the height of image is two times the height of object. The focal length of the lens (in cm) is $f$. Then $3f$ will be :
\qn{9} A variable force $\vec F=\left(x^2\hat i+xy\hat j\right)$ N is acting on a particle which moves from point $(0,0)$ to $(1\mathrm m,2\mathrm m)$ along the path $OAB$ as shown. The work done by force is $\dfrac pq$ joules, where $p$ and $q$ are co-prime integers. Find the value of $(p+q)$.
\qfig{p2_phy_s2_q09_1}
\qn{10} A ring of mass $m$ is pivoted at a point on its periphery and free to oscillates in vertical plane. A particle of same mass is attached at the lowest point of the ring then time period (in second) of small oscillations is [use $R=2$m ; $g=\pi^2$ m/s$^2$]
\qfig{p2_phy_s2_q10_1}
\raggedcolumns
\partsec{2}{CHEMISTRY}
\sect{SECTION-I}
\qn{1} \qfig{p2_chem_s1_q01_1}
If $X$ and $Y$ are functional group isomers, then what is the structure of reactant (Alkyne).
\begin{qoptions}
\opt{A}{$\mathrm{CH_3{-}C{\equiv}C{-}CH_3}$}
\opt{B}{$\mathrm{CH_3{-}CH_2{-}C{\equiv}C{-}H}$}
\opt{C}{$\mathrm{CH_2{=}CH{-}CH{=}CH_2}$}
\opt{D}{$\mathrm{CH_3{-}CH{=}C{=}CH_2}$}
\end{qoptions}
\qn{2} The major product formed in the following reaction is :
\qfig{p2_chem_s1_q02_1}
\begin{qoptions}
\opt{A}{\qfigi{p2_chem_s1_q02_A}}
\opt{B}{\qfigi{p2_chem_s1_q02_B}}
\opt{C}{\qfigi{p2_chem_s1_q02_C}}
\opt{D}{\qfigi{p2_chem_s1_q02_D}}
\end{qoptions}
\qn{3} \qfig{p2_chem_s1_q03_1}
\begin{qoptions}
\opt{A}{\qfigi{p2_chem_s1_q03_A}}
\opt{B}{\qfigi{p2_chem_s1_q03_B}}
\opt{C}{\qfigi{p2_chem_s1_q03_C}}
\opt{D}{$\mathrm{CH_2{=}CH_2}$}
\end{qoptions}
\qn{4} \qfig{p2_chem_s1_q04_1}
Final product ($B$) of the reaction sequence is :
\begin{qoptions}
\opt{A}{\qfigi{p2_chem_s1_q04_A}}
\opt{B}{\qfigi{p2_chem_s1_q04_B}}
\opt{C}{\qfigi{p2_chem_s1_q04_C}}
\opt{D}{\qfigi{p2_chem_s1_q04_D}}
\end{qoptions}
\qn{5} Which of the following alkene gives same product on reaction with $\mathrm{HBr/CCl_4}$ and $\mathrm{HBr/H_2O_2/h\nu}$
\begin{qoptions}
\opt{A}{\qfigi{p2_chem_s1_q05_A}}
\opt{B}{\qfigi{p2_chem_s1_q05_B}}
\opt{C}{\qfigi{p2_chem_s1_q05_C}}
\opt{D}{\qfigi{p2_chem_s1_q05_D}}
\end{qoptions}
\qn{6} How many moles of HI will be consumed per mole of glycerol in the following reaction.
\qfig{p2_chem_s1_q06_1}
\begin{qoptions}
\opt{A}{4}
\opt{B}{5}
\opt{C}{6}
\opt{D}{7}
\end{qoptions}
\qn{7} $\mathrm{(1)\ CH_3OH}\xrightarrow{\mathrm{Na}}A$
\qfig{p2_chem_s1_q06_2}
Product $B$ and $D$ are :
\begin{qoptions}
\opt{A}{Diastereomers}
\opt{B}{Enantiomers}
\opt{C}{Identical}
\opt{D}{Functional group isomers}
\end{qoptions}
\qn{8} Match items of \underline{List--I} with those of \underline{List--II} :
{\small\setlength{\tabcolsep}{3pt}\renewcommand{\arraystretch}{1.25}\begin{center}\begin{tabular}{|L{0.6cm}|L{3.0cm}|L{0.7cm}|L{1.5cm}|}
\hline
\multicolumn{2}{|c|}{\parbox{3.4cm}{\centering List--I\\(Property)}} & \multicolumn{2}{c|}{\parbox{2.3cm}{\centering List--II\\(Example)}}\\
\hline
(a) & Diamagnetism & (i) & $\mathrm{MnO}$\\
\hline
(b) & Ferrimagnetism & (ii) & $\mathrm{O_2}$\\
\hline
(c) & Paramagnetism & (iii) & $\mathrm{NaCl}$\\
\hline
(d) & Antiferromagnetism & (iv) & $\mathrm{Fe_3O_4}$\\
\hline
\end{tabular}\end{center}}
Choose the \textbf{most appropriate} answer from the options given below :
\begin{qoptions}
\opt{A}{(a)--(ii), (b)--(i), (c)--(iii), (d)--(iv)}
\opt{B}{(a)--(iii), (b)--(i), (c)--(ii), (d)--(iv)}
\opt{C}{(a)--(iii), (b)--(iv), (c)--(ii), (d)--(i)}
\opt{D}{(a)--(iv), (b)--(ii), (c)--(i), (d)--(iii)}
\end{qoptions}
\qn{9} Two liquids $A$ and $B$ have, $P_A^0:P_B^0=1:3$ at a certain temperature. Assume $A$ and $B$ form an ideal solution, and the ratio of moles of $A$ to $B$ in the liquid phase are $1:3$, then mole fraction of $A$ in vapour phase in equilibrium with the solution will be -
\begin{qoptions}
\opt{A}{0.1}
\opt{B}{0.2}
\opt{C}{0.5}
\opt{D}{1.0}
\end{qoptions}
\qn{10} The hydrolysis of an ester was carried out separately with 0.05 M HCl and 0.05 M $\mathrm{H_2SO_4}$. Which of the following will be true ?\\
\relax[$k$ = first order rate constant of reaction]
\begin{qoptions}
\opt{A}{$k_{\mathrm{HCl}}>k_{\mathrm{H_2SO_4}}$}
\opt{B}{$k_{\mathrm{HCl}}<k_{\mathrm{H_2SO_4}}$}
\opt{C}{$k_{\mathrm{HCl}}=k_{\mathrm{H_2SO_4}}$}
\opt{D}{$k_{\mathrm{HCl}}=2k_{\mathrm{H_2SO_4}}$}
\end{qoptions}
\qn{11} A sample of oleum is labelled as $\left(100+\dfrac{900}{89}\right)\%$. Find the moles of NaOH required to completely react with 178 gm of the oleum.
\begin{qoptions}
\opt{A}{4}
\opt{B}{6}
\opt{C}{3}
\opt{D}{7}
\end{qoptions}
\qn{12} If excess water is added into a bottle sample labelled as ``112\% oleum'' and is reacted with 1.06 gm $\mathrm{Na_2CO_3}$, then find the volume of $\mathrm{CO_2}$ evolved at 1 atm and 300 K temperature after completion of the reaction. [Assume oleum to be in excess]
\begin{qoptions}
\opt{A}{12.3 L}
\opt{B}{0.246 L}
\opt{C}{24.6 L}
\opt{D}{2.46 L}
\end{qoptions}
\qn{13} The number of Nitrogen atoms, water molecules and the total charge on the product side are respectively when following reaction is balanced ?\\
$\mathrm{Cu+HNO_3\to Cu(NO_3)_2+NO_2+H_2O}$
\begin{qoptions}
\opt{A}{6, 3, 0}
\opt{B}{4, 2, 2}
\opt{C}{4, 2, 0}
\opt{D}{3, 2, 0}
\end{qoptions}
\qn{14} Which of the following statement is \textbf{CORRECT} ?
\begin{qoptions}
\opt{A}{$\mathrm{CHF_3}$ is more acidic than $\mathrm{CHCl_3}$}
\opt{B}{Geometry around `N' is linear in $\mathrm{H_3Ge{-}NCO}$}
\opt{C}{$\mathrm{[B_3O_6]^{-3}}$ \& $\mathrm{B_3N_3H_6}$ are planar}
\opt{D}{In $\mathrm{O(SiH_3)_2}$ bond angle only affected by steric repulsion}
\end{qoptions}
\qn{15} Which of the following statement is correct?
\begin{qoptions}
\opt{A}{Octet rule is followed by N in $\mathrm{NO_2}$}
\opt{B}{$\mathrm{BF_3}$ is hypervalent species and $\mathrm{PF_5}$ is hypovalent species}
\opt{C}{Cl--O bond order in $\mathrm{ClO_3^-}$ is less than in $\mathrm{ClO_4^-}$}
\opt{D}{All of the above}
\end{qoptions}
\qn{16} Which of the following order is/are \textbf{CORRECT}?
\begin{qoptions}
\opt{A}{$\mathrm{BF_3<BCl_3<BBr_3}$ (Lewis acidic character)}
\opt{B}{$\mathrm{NO_2^+>NO_2^->NO_3^-}$ (Bond angle)}
\opt{C}{$\mathrm{{:}CH_2<{:}CF_2<{:}CCl_2}$ (Stability of Carbene)}
\opt{D}{$\mathrm{CO_3^{2-}<CO_2<CO}$ (C-O Bond length)}
\end{qoptions}
\qn{17} The nature of overlap in the bridge bond present in $\mathrm{I_2Cl_6}$ is :
\begin{qoptions}
\opt{A}{$sp^3$--$p$--$sp^3$}
\opt{B}{$sp^3d^2$--$sp^3$--$sp^3d^2$}
\opt{C}{$sp^3d$--$sp^3$--$sp^3d$}
\opt{D}{$sp^3d$--$sp^2$--$sp^3d$}
\end{qoptions}
\qn{18} According to M.O.T., which of the following molecule has maximum number of $\pi$-bond?
\begin{qoptions}
\opt{A}{$\mathrm{C_2}$}
\opt{B}{$\mathrm{O_2}$}
\opt{C}{$\mathrm{OF}$}
\opt{D}{$\mathrm{NO}$}
\end{qoptions}
\qn{19} Which of the following is not a correct order\\
(i) $\mathrm{AlCl_3<AlBr_3<AlI_3}$ (Covalent character)\\
(ii) $\mathrm{Rb_2CO_3>K_2CO_3>Na_2CO_3}$ (decomposition temperature)\\
(iii) $\mathrm{BeC_2O_4<MgC_2O_4<CaC_2O_4}$ (Solubility)\\
(iv) $\mathrm{Li_3N>Na_3N>K_3N}$ (Thermal stability)
\begin{qoptions}
\opt{A}{i, iii, iv}
\opt{B}{i, iii}
\opt{C}{only iii}
\opt{D}{iii and iv}
\end{qoptions}
\qn{20} Solid $AB$ form NaCl type structure in which $A^+$ ion is present at corner \& face center of unit cell. If this unit cell is cut in two identical parts according to given plane in the figure. These two parts of unit cell called newly generated unit cells. INCORRECT option is -
\qfig{p2_chem_s1_q20_1}
\begin{qoptions}
\opt{A}{Empirical formula of $AB$ remains unchanged in newly generated unit cell}
\opt{B}{Equal numbers of $A^+$ and $B^-$ ions are present at both edge centres and corners in newly generated unit cell}
\opt{C}{Cation $A^+$ will be present on face centres of newly generated unit cell}
\opt{D}{Newly generated unit cell will have same lattice parameters as that of tetragonal lattice system}
\end{qoptions}
\sect{SECTION-II}
\qn{1} \qfig{p2_chem_s2_q01_1}
Identify total number of compound gives immidiate turbidity with Lucas reagent.
\qn{2} \qfigi{p2_chem_s2_q02_1}\qfigi{p2_chem_s2_q02_2}\par
Total number of possible dibromo products in final product mixture ($Y$)
\columnbreak
\qn{3} \qfig{p2_chem_s2_q03_1}
Identify total number of stereoisomer's if double bond is present in cis configuration in given compound ($X$)
\qn{4} An ionic compound $\left(\overset{+}{X}\overset{-}{Y}\right)$ is doped with controlled amount of $(MY_2)$. Number of cation vacancies in 1 mol of $XY$ is found to be $6.02\times10^{21}$. If atomic mass of $X$, $Y$ and $M$ are 15,10 \& 20 a.m.u respectively :\\
(Assume volume remains unchanged)\\
(Given : $N_A=6.02\times10^{23}$) If the \% decreases in density is $x$, then the value of $10x$ is :
\qn{5} 585.5 g of 0.5 molal aqueous solution of sucrose can be cooled up to $-4^\circ$C before freezing. The amount of ice (in g) that will be separated out is \underline{\hspace{1.2cm}}.\\
\relax[Given : $K_f(\mathrm{H_2O})=1.86$ K kg mol$^{-1}$]\\
\relax[Give your answer as nearest integer].
\qn{6} A mixture of $\mathrm{H_2}$, $\mathrm{N_2}$ and $\mathrm{O_2}$ occupying 100 ml underwent reaction so as to form $\mathrm{H_2O_2}$ (l) and $\mathrm{N_2H_2}$ (g) as the only products, causing the volume of contract by 60 ml. The remaining mixture was passed through pyragallol causing a contraction of 10 ml. To the remaining mixture excess $\mathrm{H_2}$ was added and the above reaction was repeated, causing a reduction in volume of 10 ml. (No other products are formed)\\
What is the mole percent of $\mathrm{N_2}$ in the initial mixture ?
\qn{7} The dipole moment of molecule $A$--$B$ is 1.6D. If bond length of $A$--$B$ molecule is 1.625\AA, then \%ionic character is calculated to be $x\%$. Calculate $x\times0.39$ ?
\qn{8} Total number of boron atom in anionic part of borax which form $\sigma$-coordinate bond is $x$ and total number of boron atom which can form p$\pi$-p$\pi$ coordinate bond is $y$ then calculate of $x-y$
\qn{9} Find the number of non redox reaction to give their product(s) from the following :\\
(1) $\mathrm{XeF_2+H_2O\to}$\\
(2) $\mathrm{N_2O_5+H_2O\to}$\\
(3) $\mathrm{N_2O_3+H_2O\to}$\\
(4) $\mathrm{XeF_6+H_2O\to}$\\
(5) $\mathrm{XeF_4+H_2O\to}$
\qn{10} Total number of species in which highest occupied molecular orbital has gerade type symmetry.\\
$\mathrm{B_2}$, $\mathrm{N_2}$, $\mathrm{C_2}$, $\mathrm{N_2^{2+}}$, $\mathrm{O_2}$, $\mathrm{O_2^{2-}}$, $\mathrm{N_2^{2-}}$
\raggedcolumns
\partsec{3}{MATHEMATICS}
\sect{SECTION-I}
\qn{1} Let $f$ be a function defined for all $x>1$ such that $4f(x^{-1}+1)+8f(x+1)=\log_{12}x$. Then $4(f(10)+f(13)+f(17))$ is-
\begin{qoptions}
\opt{A}{1}
\opt{B}{2}
\opt{C}{3}
\opt{D}{4}
\end{qoptions}
\qn{2} Let $t_1=\left(\tan^{-1}x\right)^{\tan^{-1}x}$, $t_2=\left(\tan^{-1}x\right)^{\cot^{-1}x}$, $t_3=\left(\cot^{-1}x\right)^{\tan^{-1}x}$, $t_4=\left(\cot^{-1}x\right)^{\cot^{-1}x}$. If $x\in\left(0,\dfrac1{a^a}\right)$ where $a=2015$, then which of the following is(are) correct-
\begin{qoptions}
\opt{A}{$t_4>t_3>t_1>t_2$}
\opt{B}{$t_4>t_3$ and $t_1=t_2$}
\opt{C}{$t_1=t_2=t_3=t_4$}
\opt{D}{$t_2>t_1$ and $t_3=t_4$}
\end{qoptions}
\qn{3} If the function $f:R-\{-b\}\to R-\{1\}$ is defined by $f(x)=\dfrac{(x+a)}{(x+b)}$ ($a\ne b$) is inverse of itself, then $b+2$ is
\begin{qoptions}
\opt{A}{2}
\opt{B}{1}
\opt{C}{5}
\opt{D}{8}
\end{qoptions}
\qn{4} If $f(x)=\sqrt[5]{5-x^5}$, then \[\underbrace{f(f(f.\ldots.(1)))}_{2016\ \text{times}\ f}\] is equal to -
\begin{qoptions}
\opt{A}{$4^{1/5}$}
\opt{B}{$4^{\frac{2016}{5}}$}
\opt{C}{$4^{\frac1{10080}}$}
\opt{D}{1}
\end{qoptions}
\qn{5} If $\sin^{-1}x+\sin^{-1}y+\sin^{-1}z=\pi$, then-
\begin{qoptions}
\opt{A}{$\left(1+\sqrt{1-x^2}\right)\left(1+\sqrt{1-y^2}\right)\left(1+\sqrt{1-z^2}\right)=1$}
\opt{B}{$x\left(\sqrt{1-x^2}\right)+y\left(\sqrt{1-y^2}\right)+z\left(\sqrt{1-z^2}\right)=2xyz$}
\opt{C}{$x^2+y^2+z^2=\left(\sqrt{\left(1-x^2\right)\left(1-y^2\right)\left(1-z^2\right)}\right)$}
\opt{D}{$x+y+z=xyz$}
\end{qoptions}
\qn{6} If $\tan^{-1}(x+1)+\cot^{-1}(x-1)=\sin^{-1}\dfrac45+\cos^{-1}\dfrac35$, then $x$ has the value-
\begin{qoptions}
\opt{A}{$4\sqrt{\dfrac37}$}
\opt{B}{$4\sqrt{\dfrac73}$}
\opt{C}{$14\sqrt3$}
\opt{D}{$6\sqrt7$}
\end{qoptions}
\qn{7} Number of solution(s) of $\left(\tan^{-1}x\right)^2+\left(\cot^{-1}x\right)^2=\dfrac{5\pi^2}8$, is-
\begin{qoptions}
\opt{A}{1}
\opt{B}{2}
\opt{C}{3}
\opt{D}{4}
\end{qoptions}
\qn{8} Value of $\displaystyle\lim_{x\to0}\dfrac{\ln\left(1+3x^2\right)\sin\left(\tan^{-1}x\right)}{(3^x-1)(1-\cos2x)}\left[\dfrac{\tan x}{x}\right]$ is (where $[.]$ represents greatest integer function) -
\begin{qoptions}
\opt{A}{$\dfrac32$}
\opt{B}{$\dfrac12$}
\opt{C}{$\dfrac3{\ln3}$}
\opt{D}{$\dfrac3{2\ln3}$}
\end{qoptions}
\qn{9} If $\displaystyle\lim_{x\to0}\dfrac{\alpha\sin x-\beta x\cos x}{\tan x-\sin x}=1$, then $(\alpha+\beta)$ is equal to-
\begin{qoptions}
\opt{A}{1}
\opt{B}{$-1$}
\opt{C}{2}
\opt{D}{3}
\end{qoptions}
\qn{10} Value of $\displaystyle\lim_{x\to0}\dfrac{(\sin2x)\left(\tan x-\tan^{-1}x\right)}{\left(e^{\cos3x}-e\right)(\ln\cos x)}$ is equal to-
\begin{qoptions}
\opt{A}{$\dfrac{16}{27e}$}
\opt{B}{$\dfrac{16}{27}$}
\opt{C}{$\dfrac4{9e}$}
\opt{D}{$\dfrac49$}
\end{qoptions}
\qn{11} Value of $\displaystyle\lim_{x\to0}\dfrac{\sqrt{\cos2x}-\sqrt{1+x\sin x}}{x\tan\frac x2}$ is -
\begin{qoptions}
\opt{A}{3}
\opt{B}{$-3$}
\opt{C}{6}
\opt{D}{$-6$}
\end{qoptions}
\qn{12} The value of $\displaystyle\lim_{x\to1}\dfrac{x^3-1}{(x^2-3x+2)-\sin(x-1)}$ is-
\begin{qoptions}
\opt{A}{$-\dfrac32$}
\opt{B}{1}
\opt{C}{2}
\opt{D}{3}
\end{qoptions}
\qn{13} Value of $\displaystyle\lim_{x\to0}\dfrac{\sin(2\pi\sec x)}{x^2}$ is-
\begin{qoptions}
\opt{A}{$\dfrac\pi4$}
\opt{B}{$\dfrac\pi2$}
\opt{C}{$\pi$}
\opt{D}{$2\pi$}
\end{qoptions}
\qn{14} Range of $f(x)=k\sin x+\cos x$ is $[-5,5]$, then total number of integers between the two possible value of `$k$' is-
\begin{qoptions}
\opt{A}{11}
\opt{B}{10}
\opt{C}{9}
\opt{D}{5}
\end{qoptions}
\qn{15} If $\sin A+\cos B=a$ and $\sin B+\cos A=b$, then $\sin(A+B)$ is equal to-
\begin{qoptions}
\opt{A}{$\dfrac{a^2+b^2}2$}
\opt{B}{$\dfrac{a^2-b^2+2}2$}
\opt{C}{$\dfrac{a^2+b^2-2}2$}
\opt{D}{$\dfrac{a^2+b^2+2}2$}
\end{qoptions}
\qn{16} The value of $\cos\dfrac\pi5\cos\dfrac{2\pi}5\cos\dfrac{4\pi}5\cos\dfrac{8\pi}5$ is equal to-
\begin{qoptions}
\opt{A}{$\dfrac1{16}$}
\opt{B}{$\dfrac14$}
\opt{C}{$-\dfrac18$}
\opt{D}{$-\dfrac1{16}$}
\end{qoptions}
\qn{17} Number of solutions of the equation $\tan\theta+\tan2\theta+\tan3\theta=\tan\theta\tan2\theta\tan3\theta$ in $(0,\pi)$ is-
\begin{qoptions}
\opt{A}{5}
\opt{B}{3}
\opt{C}{2}
\opt{D}{zero}
\end{qoptions}
\qn{18} The number of solutions of $2^{\sin|x|}=4^{|\cos x|}$ in $[-\pi,\pi]$ is -
\begin{qoptions}
\opt{A}{0}
\opt{B}{2}
\opt{C}{4}
\opt{D}{6}
\end{qoptions}
\qn{19} The general solution of the equation $\sin x-\cos x=\min\left\{\dfrac1{\sqrt2},x^2-6x+10\right\}$ where $n\in I$, is-
\begin{qoptions}
\opt{A}{$\dfrac\pi4+n\pi+(-1)^n\dfrac\pi6$}
\opt{B}{$n\pi+(-1)^n\dfrac\pi6-\dfrac\pi4$}
\opt{C}{$2n\pi$}
\opt{D}{$2n\pi+\dfrac\pi2$}
\end{qoptions}
\qn{20} As shown in the figure, the area of the trapezium $ABCD$ is $\dfrac{81}2$, where $a=2b$, then $c$ is equal to -
\qfig{p2_math_s1_q20_1}
\begin{qoptions}
\opt{A}{$\dfrac92$}
\opt{B}{$3\sqrt5$}
\opt{C}{$\sqrt{30}$}
\opt{D}{none of these}
\end{qoptions}
\sect{SECTION-II}
\qn{1} Let {\small\[f(x)=\cot\left(\dfrac{\cos^{-1}\sin\left(\cos^{-1}x\right)+\sin^{-1}\left(\cos\left(\sin^{-1}x\right)\right)}{2}\right),\]} then find the value of $\displaystyle\sum_{x=1}^{7}f\left(\dfrac x8\right)$.
\qn{2} The equation $\sin\cot^{-1}\cos\tan^{-1}x=a$ has real solution iff complete range of $a$ is $[p,q]$. Value of $\left[q+\dfrac1p\right]$ is (where $[.]$ denotes greatest integer function)
\qn{3} If $\sin^{-1}\dfrac1{\sqrt{10}}+\sin^{-1}\dfrac2{\sqrt{85}}+\sin^{-1}\left(\dfrac4{\sqrt{1105}}\right)+\ldots\ldots\infty=\alpha$, then value of $\tan\alpha$ is
\qn{4} If $\tan^{-1}x+\tan^{-1}y=\cot^{-1}z$ is such that $x+y+z=\sqrt3$, then value of $\dfrac1{xyz}$ can be expressed as $k\sqrt k$, where $k$ is
\qn{5} The whole domain of $y=f(x)$ defined implicitly by the equation $3^y+2^{x^4}=2^{4x^2-1}$ is $(a,b)\cup(c,d)$, then $a+b+c+d$ is
\qn{6} Value of $\displaystyle\lim_{x\to0}\dfrac{x^2(x-\tan x)}{(x-\sin x)}\cdot\dfrac{\left(\sin^{-1}x\right)^2}{x^2-\left(\sin^{-1}x\right)^2}$ is
\qn{7} The value of $4e\left(\displaystyle\lim_{x\to0}\left(\dfrac{\sin x}{x}\right)^{\frac{\sin x}{x-\sin x}}+\lim_{x\to1}(x)^{\frac1{1-x}}\right)$ is
\qn{8} If in a triangle $ABC$ with usual notation, $a(a-c)=b^2-c^2$, then $1572(\sin^2(\pi\cos^2B))$ is equal to
\qn{9} If $r_1=4$, $r_2=6$, $r_3=8$ for $\Delta ABC$ and $r=\dfrac pq$ (where $p$ \& $q$ are relatively prime numbers), then value of $(2q-p)$ is (usual notation is followed)
\qn{10} If in a $\Delta ABC$ with usual notations, $3a+3b=4c$, then $\cot\dfrac A2\cdot\cot\dfrac B2$ is
\end{multicols}
\end{document}
